\section{Complete fixed-kernel \texorpdfstring{$D_3$}{D3} Gauge-index Classification and Conditional Reduction}
\label{app:d3-coefficient-rigidity}

We set
\begin{equation}
 \rho=\frac{\epsilon_V^2}{N_{\rm gen}}>0.
 \label{eq:pr35-d3-rho}
\end{equation}
This appendix classifies every real symmetric gauge-index coefficient matrix multiplying one fixed Lorentz/momentum kernel at the declared scalar order.  It
does not exhaust independent Lorentz structures, CP-odd or complex Hermitian
kernels, higher derivatives, nonlocal operators, or additional fields; those
would require a separate operator-basis theorem.  Nor does it claim uniqueness
of equation syntax. Algebraically equivalent equations necessarily describe
the same tensor.

Let
\begin{equation}
 V_{\rm EW}=\operatorname{span}_{\mathbb R}
 \{W^1,W^2,W^3,B\}=\mathcal C\oplus\mathcal N,
 \qquad \mathcal C=\operatorname{span}\{W^1,W^2\}.
\end{equation}
The canonical real inner product makes $(W^1,W^2,W^3,B)$ orthonormal.  For a
unit vector $v$, let $P_v=v\otimes v$ be its orthogonal projector.  Then
\begin{equation}
 \|P_v\|_F^2=\operatorname{Tr}(P_v^2)=1,
 \qquad P_{\rm ch}=P_{W^1}+P_{W^2}.
 \label{eq:pr35-d3-projector-normalization}
\end{equation}
PR-II fixes the physical neutral basis \cite{PRII}
\begin{equation}
 Z=cW^3-sB,\qquad A=sW^3+cB,
 \qquad c=\cos\theta_W,\quad s=\sin\theta_W.
 \label{eq:pr35-d3-neutral-basis}
\end{equation}

\subsection{Exhaustion Under Residual Symmetry}

In the ordered basis $(W^1,W^2,Z,A)$, residual $U(1)_{\rm em}$ acts as
\begin{equation}
 R(\alpha)=\operatorname{diag}(R_2(\alpha),1,1),\qquad
 J=R'(0)=\operatorname{diag}(J_2,0,0),\qquad
 J_2=\begin{pmatrix}0&-1\\1&0\end{pmatrix}.
\end{equation}

\begin{theorem}[Complete residual-$U(1)$ classification]
\label{thm:pr35-d3-u1-classification}
For $D\in\operatorname{Sym}(V_{\rm EW})$, the identity
\begin{equation}
 R(\alpha)^TDR(\alpha)=D\qquad\text{for every }\alpha\in\mathbb R
 \label{eq:pr35-d3-u1-invariance}
\end{equation}
holds if and only if
\begin{equation}
 D=aP_{\rm ch}+zP_Z+x(Z\otimes A+A\otimes Z)+\gamma P_A
 \label{eq:pr35-d3-complete-invariant-family}
\end{equation}
for arbitrary $a,z,x,\gamma\in\mathbb R$.
\end{theorem}

\begin{proof}
Write the symmetric block matrix
\begin{equation}
 D=\begin{pmatrix}C&L\\L^T&N\end{pmatrix},
 \qquad C,N\in\operatorname{Sym}_2(\mathbb R).
\end{equation}
Differentiating Eq.~\eqref{eq:pr35-d3-u1-invariance} at $\alpha=0$ gives
$[D,J]=0$, hence
\begin{equation}
 CJ_2=J_2C,\qquad J_2L=0,\qquad L^TJ_2=0.
\end{equation}
Because $J_2$ is invertible, $L=0$.  If
$C=\left(\begin{smallmatrix}u&v\\v&w\end{smallmatrix}\right)$, the first
equation gives $v=0$ and $u=w$, so $C=aI_2$.  The residual group acts trivially on the neutral plane, leaving $N$ as an arbitrary real symmetric $2\times2$ matrix.  Expanding $N$ in the $(Z,A)$ basis gives exactly
Eq.~\eqref{eq:pr35-d3-complete-invariant-family}.  Direct substitution proves
the converse.
\end{proof}

The accompanying exact rational row reduction starts with all ten independent
components of a symmetric $4\times4$ matrix.  Equation
\eqref{eq:pr35-d3-u1-invariance} has rank six and therefore the four-dimensional
solution space in Theorem~\ref{thm:pr35-d3-u1-classification}.  Imposing an
unchanged charged sector,
\begin{equation}
 P_{\rm ch}DP_{\rm ch}=0,
 \label{eq:pr35-d3-charged-unchanged}
\end{equation}
sets $a=0$; the combined rank is seven and the solution dimension is three. If $D_3$ is a zero-momentum mass-kernel correction, photon preservation adds
\begin{equation}
 DA=0.
 \label{eq:pr35-d3-photon-null}
\end{equation}
Equation~\eqref{eq:pr35-d3-complete-invariant-family} then gives
$x=\gamma=0$.  The combined rank is nine, and the complete structural family
is therefore
\begin{equation}
 D(k)=\rho kP_Z,\qquad k\in\mathbb R.
 \label{eq:pr35-d3-all-real-k-family}
\end{equation}
Every real $k$ satisfies residual symmetry, charged-sector preservation, and
photon nullity.  These conditions fix a one-dimensional subspace, not its
scalar.  For a
momentum-dependent polarization, ordinary Lorentz transversality must not be
silently replaced by Eq.~\eqref{eq:pr35-d3-photon-null}; the displayed result
uses the stronger mass-kernel condition relevant to the PR-III scalar mass
correction.

\begin{theorem}[Frozen-reduct expansions]
\label{thm:pr35-d3-common-reduct}
Let $\mathcal L_0$ be the target-free language containing P1 through P7 and the frozen background objects used above---$\rho$, the canonical metric, charge grading, and physical $Z,A$ basis---but no $D_3$ tensor symbol.  Let $\mathsf T_0$ be the corresponding public reduct, excluding the candidate equation $D_{3,Z}=2\rho$ and downstream diagnostic comparisons.  Extend the language by
a real symmetric symbol $D_3$, and set
\begin{equation}
 \mathsf T_{\rm struct}=\mathsf T_0\cup
 \left\{\begin{array}{l}
 R(\alpha)^TD_3R(\alpha)=D_3\quad(\forall\alpha\in\mathbb R),\\
 P_{\rm ch}D_3P_{\rm ch}=0,\quad D_3A=0
 \end{array}\right\}.
 \label{eq:pr35-d3-structural-theory}
\end{equation}
The theory $\mathsf T_0$ is nonempty.  For every $k\in\mathbb R$, the
expansion of the explicit model $\mathfrak M_0$ constructed below,
\begin{equation}
 \mathfrak M_k=(\mathfrak M_0,D_3=\rho kP_Z)
 \label{eq:pr35-d3-model-expansion}
\end{equation}
models $\mathsf T_{\rm struct}$.  Moreover, $\mathfrak M_k$ and
$\mathfrak M_{k'}$ are inequivalent under normalized orthogonal basis changes
whenever $k\ne k'$.
\end{theorem}

\begin{proof}
Start with the explicit P1 through P7 common reduct constructed in Appendix~\ref{app:axiom-by-axiom}.  Adjoin the independent target-free
background $V_{\rm EW}=\mathbb R^4$ with orthonormal basis $(W^1,W^2,W^3,B)$; choose $\rho>0$ and $c^2+s^2=1$; and define $Z,A$ and $R(\alpha)$ by Eqs.~\eqref{eq:pr35-d3-neutral-basis} and the residual action above.  These definitions alter no old-language symbol and contain no $D_3$ equation, so they give an explicit $\mathfrak M_0\models\mathsf T_0$. Changing the interpretation of the new $D_3$ symbol cannot alter any $\mathcal L_0$ sentence, so every expansion preserves all P1--P7 and target-free background statements true in $\mathfrak M_0$.  Equation \eqref{eq:pr35-d3-all-real-k-family} verifies each added structural sentence.
Finally,
\begin{equation}
 \operatorname{spec}_{\rm mult}(D_3/\rho)=\{k,0,0,0\}.
\end{equation}
Normalized orthogonal similarity preserves this eigenvalue multiset, proving
inequivalence when $k\ne k'$.
\end{proof}

This is the axiom-by-axiom common-reduct proof of an uncountable structural
model class.  If Eq.~\eqref{eq:pr35-d3-photon-null} is not a derived
mass-kernel condition, that class is larger still.

\subsection{Canonical Adjoint Trace}

PR-II fixes \cite{PRII}
\begin{equation}
 T_a=\frac{\sigma_a}{2},\qquad
 \operatorname{Tr}(T_aT_b)=\frac12\delta_{ab},\qquad
 [T_a,T_b]=i\epsilon_{abc}T_c,
 \label{eq:pr35-d3-prii-normalization}
\end{equation}
and places the weak projection modes in the adjoint $j=1$ representation.
On $\mathcal H_{\rm ch}=\operatorname{span}\{T^+,T^-\}$,
$[Q,T^+]=+T^+$ and $[Q,T^-]=-T^-$.  Let $e_3=(0,0,1)^T$ and
$P_3=e_3\otimes e_3$.  Define
\begin{equation}
 (F_a)_{bc}=-i\epsilon_{abc},\qquad
 P_{\rm ch}^{\rm ad}=\operatorname{diag}(1,1,0),\qquad
 K_{ab}=\operatorname{Tr}_{\rm ad}
 \!\left(P_{\rm ch}^{\rm ad}F_aP_{\rm ch}^{\rm ad}F_b\right).
 \label{eq:pr35-d3-projected-tensor}
\end{equation}
Here $P_3$ denotes the adjoint third-axis projector and is unrelated to the
base postulate carrying the label P3.

\begin{lemma}[Exact projected-adjoint tensor]
\label{lem:pr35-d3-adjoint-trace}
With the frozen normalization in Eq.~\eqref{eq:pr35-d3-prii-normalization},
\begin{equation}
 K_{ab}=2\delta_{a3}\delta_{b3}=2(P_3)_{ab}.
 \label{eq:pr35-d3-k-equals-two-p3}
\end{equation}
\end{lemma}

\begin{proof}
$P_{\rm ch}^{\rm ad}F_1P_{\rm ch}^{\rm ad}$ and
$P_{\rm ch}^{\rm ad}F_2P_{\rm ch}^{\rm ad}$ vanish, whereas
\begin{equation}
 P_{\rm ch}^{\rm ad}F_3P_{\rm ch}^{\rm ad}
 =\begin{pmatrix}0&-i&0\\i&0&0\\0&0&0\end{pmatrix},
 \qquad
 \operatorname{Tr}\!\left[
 (P_{\rm ch}^{\rm ad}F_3P_{\rm ch}^{\rm ad})^2\right]=2.
\end{equation}
Thus only $K_{33}$ is nonzero and it equals two.  Equivalently, the charged
eigenvalues of $\operatorname{ad}_Q$ are $(+1,-1)$, so their squared trace is
$1^2+(-1)^2=2$; or directly
$\sum_{b,c}\epsilon_{3bc}^2=\epsilon_{312}^2+\epsilon_{321}^2=2$.
\end{proof}

The charged scalar trace is invariant under normalized transformations that
preserve $Q$ and the charged subspace.  The displayed component tensor is
covariant relative to the distinguished neutral axis.  The lemma fixes a
canonical representation trace; it does not fix the dynamical scalar
multiplying an effective operator.

\subsection{Conditional Reduction to the Canonical Tensor}

Let
\begin{equation}
 \mathcal E:\operatorname{span}_{\mathbb R}\{P_3\}
 \longrightarrow\operatorname{span}_{\mathbb R}\{P_Z\}
\end{equation}
be a specified normalized real-linear isometry.  The needed hypotheses are:
\begin{enumerate}[label=\textnormal{H\arabic*.},leftmargin=2.8em]
 \item the complete order-$\rho$ neutral operator basis contains exactly one
 contribution;
 \item the full determinant Hessian, including Taylor, vector, ghost,
 Goldstone \cite{PeskinSchroeder_1995}, matter, spectral-kernel, sign, coupling, and symmetry factors, is
 exactly $D=\rho\mathcal E(K)$ with coefficient $+1$;
 \item $\mathcal E$ is real-linear and $\mathcal E(P_3)=P_Z$, as a normalized
 neutral-line map rather than an unannounced ordinary $W^3$ component
 projection;
 \item generator, trace, field, kinetic, $\epsilon_V$, and generation
 normalizations are frozen; and
 \item Eqs.~\eqref{eq:pr35-d3-u1-invariance}--
 \eqref{eq:pr35-d3-photon-null} hold.
\end{enumerate}

\begin{theorem}[Conditional canonical $D_3$ reduction]
\label{thm:pr35-d3-conditional-uniqueness}
Under H1--H5,
\begin{equation}
 D=2\rho P_Z,\qquad
 D_{3,W}=0,\qquad
 D_{3,Z}=\frac{2}{N_{\rm gen}}\epsilon_V^2,
 \label{eq:pr35-d3-unique-canonical-result}
\end{equation}
so $k=2$ is the unique coefficient in this operator class.
\end{theorem}

\begin{proof}
Lemma~\ref{lem:pr35-d3-adjoint-trace}, H2, and H3 give
\begin{equation}
 D=\rho\mathcal E(K)=2\rho\mathcal E(P_3)=2\rho P_Z.
\end{equation}
H1 excludes a second same-order tensor, and H4 excludes absorbing a multiplier
into the fields or $\epsilon_V$.
\end{proof}

This theorem uses no experimental residual and is not a best-fit selection.
It is deliberately conditional: H2 is precisely the missing dynamical
unit-prefactor Hessian identification, not a result obtained inside this
reduction.  Once H2 is independently derived from the determinant, the
remaining reduction to two is exact.

\subsection{Necessary-and-sufficient Candidate Characterization}

For arbitrary $D\in\operatorname{Sym}(V_{\rm EW})$, define
\begin{align}
 E_{\rm sym}(D)&=\sup_{\alpha}
 \|R(\alpha)^TDR(\alpha)-D\|_F,\nonumber\\
 E_{\rm ch}(D)&=\|P_{\rm ch}DP_{\rm ch}\|_F,\nonumber\\
 E_{\gamma}(D)&=\|DA\|_2,\nonumber\\
 E_{\rm norm}(D)&=|Z^TDZ-2\rho|.
 \label{eq:pr35-d3-four-defects}
\end{align}

\begin{theorem}[Complete candidate characterization]
\label{thm:pr35-d3-four-defect-certificate}
All four defects in Eq.~\eqref{eq:pr35-d3-four-defects} vanish if and only if
$D=2\rho P_Z$.
\end{theorem}

\begin{proof}
$E_{\rm sym}=0$ gives the complete family in
Theorem~\ref{thm:pr35-d3-u1-classification}.  Then $E_{\rm ch}=0$ sets $a=0$,
and $E_\gamma=0$ sets $x=\gamma=0$, leaving $D=zP_Z$.
$E_{\rm norm}=0$ finally gives $z=2\rho$.  Direct substitution proves the
converse.
\end{proof}

Thus every other real symmetric tensor either violates a structural condition
or differs from the stipulated candidate normalization.
On the structural family \eqref{eq:pr35-d3-all-real-k-family},
\begin{equation}
 D(k)-D(2)=\rho(k-2)P_Z,\qquad
 \|D(k)-D(2)\|_F^2=\rho^2(k-2)^2,
 \label{eq:pr35-d3-complete-real-defect}
\end{equation}
whose sole real zero is $k=2$.  The logical boundary is decisive:
$E_{\rm norm}=0$ must be derived from the projection determinant.  Merely
declaring it assumes the desired coefficient.  This theorem is therefore an
exact candidate characterization, not by itself a postulate-derived physical
exclusion.

\subsection{Exact Countermodels and Electroweak Mixing No-Go}

Without H2, the complete scalar counterfamily
\begin{equation}
 D_\lambda=\rho\lambda\mathcal E(K)=2\lambda\rho P_Z,
 \qquad \lambda\in\mathbb R,
 \label{eq:pr35-d3-lambda-counterfamily}
\end{equation}
passes the first three defects.  Fixed rational choices require no fitted continuous parameter: $\lambda=1/2$ gives $k=1$, and $\lambda=3/2$ gives $k=3$.  They witness failure of entailment; they are not proposed replacements for the published model. The three linear structural gates in Theorem~\ref{thm:pr35-d3-common-reduct} permit every real $k$.  An independent strict positivity condition on the neutral correction would restrict the family to $k>0$, while positive semidefiniteness would give $k\geq0$.  Either case remains uncountable and retains the fixed alternatives $k=1$, $2\cos^2\theta_W$, $2$, and $3$; neither selects two. Checking only the stated $ZZ$ scalar is still weaker.  Every tensor
\begin{equation}
 \rho\{2P_Z+x(Z\otimes A+A\otimes Z)+\gamma P_A\}
\end{equation}
has the same $ZZ$ entry for arbitrary $x,\gamma$.  Photon nullity eliminates these two directions, but leaves the scalar family
\eqref{eq:pr35-d3-lambda-counterfamily}.

The ordinary electroweak rotation creates a separate normalization issue.  In
the $(Z,A)$ basis,
\begin{equation}
 2P_{W^3}=2\begin{pmatrix}c^2&cs\\cs&s^2\end{pmatrix}.
 \label{eq:pr35-d3-literal-w3-embedding}
\end{equation}
Its $ZZ$ pullback is $2c^2$, but it has $ZA$ mixing and a photon component.
A photon-safe tensor with the same scalar is $2c^2P_Z$, which is a different
operator and still fails the canonical normalization identity.

\begin{proposition}[Mixing normalization no-go]
\label{prop:pr35-d3-mixing-nogo}
Assume $D$ is real symmetric, $0<c<1$, and $DA=0$.  It is impossible to have
both $D_{W^3W^3}=2\rho$ and $D_{ZZ}=2\rho$.
\end{proposition}

\begin{proof}
Since $W^3=cZ+sA$, symmetry and $DA=0$ give
\begin{equation}
 D_{W^3W^3}=(cZ+sA)^TD(cZ+sA)=c^2D_{ZZ}.
\end{equation}
The two proposed nonzero equalities therefore imply $c^2=1$, a contradiction.
\end{proof}

Consequently three inequivalent maps give, respectively,
\begin{equation}
 \begin{array}{lll}
 P_3\mapsto P_Z:&D=2\rho P_Z,&k=2,\\
 \text{component-preserving photon completion}:&
 D=(2\rho/c^2)P_Z,&k=2/c^2,\\
 \text{literal }W^3\text{ embedding}:&D=2\rho P_{W^3},&
 D_{ZZ}/\rho=2c^2.
 \end{array}
 \label{eq:pr35-d3-three-neutral-maps}
\end{equation}
The public chain must derive which map is meant and where the weak-mixing
factor is normalized.  The literal $2P_{W^3}$ tensor and the photon-safe
$2c^2P_Z$ tensor are distinct cases: the first fails photon nullity, while the
second passes it but changes the normalization.

\subsection{Locked Diagnostics do not Isolate the Coefficient}

The released arithmetic is
\begin{equation}
 M_Z(k)=M_{Z,C1}\sqrt{1+\Pi_{Z,C2}+\rho k}.
 \label{eq:pr35-d3-mz-k}
\end{equation}
Monotonicity of the positive square root gives the exact released strict
one-standard-deviation open interval
\begin{equation}
 0.768434787688685769\ldots< k <
 3.145092343548119919\ldots .
 \label{eq:pr35-d3-diagnostic-interval}
\end{equation}
It is uncountable and its integer members are exactly $\{1,2,3\}$.  Its
endpoints have unit diagnostic distance and fail the released strict
acceptance gate.  The exact
diagnostic central-value solution is
$k=1.956749882341106\ldots$, not two.  Representative results are
\begin{table}[H]
\centering
\caption{Parameter-free $D_3$ normalizations under the locked PR-III mass diagnostic.}
\label{tab:pr35-d3-countermodel-diagnostics}
\begin{tabular}{@{}lrr@{}}
\toprule
$k$ & $M_Z(k)$ (GeV) & absolute diagnostic distance ($\sigma$)\\
\midrule
$1$ & $91.1859092277010$ & $0.80513$\\
$2\cos^2\theta_W=1.53728847\ldots$ & $91.1868587299446$ & $0.35299$\\
$2$ & $91.1876764310439$ & $0.03640$\\
$3$ & $91.1894436001394$ & $0.87790$\\
\bottomrule
\end{tabular}
\end{table}
$M_W$ and $G_F$ are independent of $k$ in the released generator.  Thus the
locked diagnostics do not prove uniqueness and do not justify describing two
as a best-fit selection.  The PR-I quasar diagnostic is not connected to this
PR-III coefficient by a released operator or propagation equation.

\subsection{Public-source Ancestry and Exact Verdict}

The released chain establishes the representation data and the stated
candidate, but not the dynamical link between them.  PR-III gives the
reference-normalized determinant
\begin{equation}
 \Delta\Gamma_{\rm PR}^{(1)}[B]
 =\frac12\operatorname{STr}_{\rm PR}'
 \log[H_{\rm PR}[B]H_{\rm PR}[0]^{-1}],
\end{equation}
including vector, ghost, scalar, and fermion blocks \cite{PeskinSchroeder_1995, PRIII}.  Its electromagnetic charged-mode inventory is a $Q_i^2$ extraction for the photon kinetic term, not a derivation of the neutral $D_3$ mass Hessian.  The manuscript then states Eq.~\ref{eq:pr35-pr3-structural-tensor}; the public generator assigns
\begin{equation}
 d_{3z} = \frac{2}{n_{\text{gen}}} \epsilon_v^2
\end{equation}
and marks it as a candidate rather than an exact theorem.  PR-III separately
lists derivation of the full charged and neutral self-energy kernels as future
work and explicitly distinguishes no-fit agreement from uniqueness. The missing noncircular edge is
\begin{equation}
 \frac12\operatorname{STr}'\log(H/H_0)
 \ \Longrightarrow\ \text{complete order-}\rho\text{ neutral Hessian}
 \ \Longrightarrow\ \rho\mathcal E(K)
 \label{eq:pr35-d3-missing-source-edge}
\end{equation}
with unit prefactor, no other same-order term, and
$\mathcal E(P_3)=P_Z$.  Until all H1--H5 ancestry is derived, the result is a
complete conditional reduction theorem and a complete unconditional
nonuniqueness theorem for the target-free public reduct plus the stated
structural gates.  Theorem~\ref{thm:pr35-d3-common-reduct} supplies the formal
axiom-by-axiom expansion proof.  In particular:
\begin{itemize}
 \item $K=2P_3$ is proved exactly from the frozen representation;
 \item every $k\ne2$ fails the canonical unit-prefactor identity;
 \item every tensor other than $2\rho P_Z$ fails a defining characterization
 defect, while the normalization defect remains explicitly underived physics;
 but
 \item even after imposing the three stipulated structural gates, the exact family
 $D(k)=\rho kP_Z$, and the released diagnostics permit a continuum of it.
\end{itemize}
Claiming unconditional P1--P7 uniqueness before deriving Eq.~\eqref{eq:pr35-d3-missing-source-edge} would contradict the explicit counterfamily \eqref{eq:pr35-d3-lambda-counterfamily}. The released-manuscript anchors are the PR-II generator-algebra, neutral-mixing, and charged-ladder sections, and the PR-III charged-mode-supertrace and neutral $D_3$ equations.  Their stable labels and exact source-line ranges are listed in the accompanying public-source map.
